\section{Introduction}\label{sec:intro}

Housing turnover is usually measured through market sales. Families also
transfer homes by gift, succession, and other non-market conveyances. These
deeds can change legal title without generating a market price or changing
who occupies the property. Between 2007 and 2023, our classification of
CoreLogic residential records identifies nearly 20 million family-transfer
deeds, one for every 3.6 market sales. Understanding this channel requires
measuring both its scale and the institutions that shape it.

This paper has three parts. First, it establishes national facts about
within-family transfers using 144 million residential parcel-date events.
Family deeds represent a substantial and persistent flow alongside market
sales. Their ratio to market sales is highest in California, where the tax
code historically allowed heirs to inherit a property's assessment as well
as the property. Their share rises across middle-aged housing vintages,
although it is not monotone over all ages. Many properties remain unsold
long after a family transfer. These patterns describe a channel of housing
ownership; none by itself identifies inefficient use or a causal tax effect.

Second, the paper examines how tax incentives shape recorded family
transfers using California's Proposition~19. Under Propositions~58 and 193,
eligible intergenerational transfers could preserve the parent's assessed
value. Proposition~19 narrowed this exclusion in February 2021, retaining
conditional protection for a principal residence occupied by the heir and
for family farms while removing it for other qualifying transfers
\citep{boe2025prop19}. The LAO estimated that 60,000--80,000 properties
annually had received inheritance exclusions under the earlier regime
\citep{lao2017inheritance}. That is an administrative property count, not a
tax-revenue estimate or a count of inefficiently allocated homes.

California family deeds surged 31 percent above the transparent
counterfactual in the four months before the deadline. They then fell to a
persistently lower level. A synthetic control estimates a 33 percent
2022--2023 gap; adjusting for the population aged 65 and over yields 26
percent. These are separate specifications, not the endpoints of a
confidence interval. A DiD adjustment for California's market-sale decline
implies roughly \WelfareFlow\ missing deeds annually when scaled by the
2018--2019 baseline. Inference varies across designs: the raw SCM
post/pre-MSPE rank has $p=0.039$, while the net-DiD state-placebo rank has
$p=0.176$. Conditional sale hazards and estate/trust-seller market sales do
not establish an accompanying release of homes through 2023. Deferred
succession is consistent with the evidence but is not identified by those
null results.

Third, a simple framework quantifies the potential allocation gains from
removing the inherited-assessment advantage. The object is the distortion
created by the tax wedge, not the existence of non-market transfers. Family
attachment, insurance, avoided transaction costs, and efficient rental use
can make retention socially valuable. Inefficient retention arises only
when the tax advantage induces a family to forgo an alternative use with
greater net surplus. Recorded deeds do not directly identify that margin.

We therefore combine the estimated deed response with explicit assumptions
about the conversion to distinct allocation-changing properties, the tax
wedge, succession delay, and persistence. California parcel episodes supply
a descriptive persistence proxy, and an LAO example for long-held Los
Angeles homes supplies an external wedge benchmark. Under a uniform
valuation-gap approximation, the annual gain per marginal property is half
the wedge. This sufficient-statistics-inspired calculation retains the
model assumptions needed for welfare interpretation
\citep{chetty2009sufficient}. The resulting tables show annual gains,
property-years, and discounted gains, including zero when deeds change but
allocations do not. They are conditional policy calculations, not estimates
of realized or total Proposition~19 welfare.

The paper connects the economics of bequests and inter-vivos transfers
\citep{bernheimshleifersummers1985bequest,mcgarry1999intervivos,kopczuk2007bequest,piketty2011inheritance}
to an asset-level record of housing ownership. Research on intergenerational
wealth transmission \citep{charleshurst2003wealth} motivates the importance
of the asset, while evidence on heirs' property documents a distinct source
of title-related immobility \citep{dobbsgaither2023heirs,burnettwm2025heirs}.
Recent work studies intergenerational housing-capital persistence using
linked Census, property, and income-tax records
\citep{binderrischvoorheis2026housing}. Our deed-based object is the
conveyance of particular assets, rather than the persistence of children's
housing holdings. Closely related work by
\citet{bakerkimbiswas2026inherited} uses parent-child exclusion claims in
Los Angeles and San Francisco to study Proposition~19 retiming and
subsequent sales. The present national measurement and California-wide
comparisons complement that evidence; our recorder-based family proxy
covers a broader population than confirmed parent-child exclusion claims.
The nearest descriptive antecedents also include industry tabulations
\citep{delventhal2026inherited} and the LAO's assessment-roll accounting.

The analysis also extends the property-tax and housing lock-in literature
\citep{wasiwhite2005prop13,ferreira2010takeit,fonsecaliu2024lockin,covenetal2024proptax}
to intergenerational title transfers. The transaction-tax literature
\citep{kopczukmunroe2015mansion,slemrodwebershan2017dc,bestkleven2018notches,hilberlyytikainen2017transfer}
provides a comparison for timing and volume responses, although a deadline
response alone does not reveal the valuation gaps required for welfare.
Misallocation concerns, as in \citet{glaeserluttmer2003rentcontrol}, require
evidence or assumptions about forgone surplus. Older housing can serve an
important filtering role \citep{rosenthal2014filtering}, but a family-owned
rental may already provide those services to a tenant.

Section~\ref{sec:data} describes the data and classification.
Section~\ref{sec:facts} presents the four national facts.
Section~\ref{sec:prop19} studies the policy response, and
Section~\ref{sec:welfare} develops the framework and conditional welfare
calibration. Section~\ref{sec:conclusion} concludes.
